\section{What Determines the Effective History Span}
\label{sec:why}

\subsection{A single-axis theory fails}
\label{sec:bocpd}

The standard account of non-stationarity says that old data becomes useless when the regime changes. Formalizing it with Bayesian online changepoint detection \citep{bocpd} gives a sharp hypothesis: the useful history is the time since the last change, and the effective span should track the posterior expected run length $\mathbb{E}[r_t]$ of the regime process. We compute $\mathbb{E}[r_t]$ per dataset with a Gaussian unknown-mean-and-variance observation model across a hazard sweep (Appendix \ref{app:bocpd}). The hypothesis fails. The two datasets with the \emph{shortest} expected run lengths at hazard $1/50$, traffic (48 steps) and electricity (53), are precisely the ones whose optimal lookback is the \emph{longest} (2880), yielding a negative rank correlation on one model family (Spearman $\rho = -0.67$ at hazard $1/50$; we report this as pattern evidence rather than significance: the exact permutation $p$ is 0.08 and does not survive multiple-comparison control, while the negative direction persists across the hazard sweep and flips sign at the most extreme hazard). Two robustness checks close the loop (Appendix \ref{app:bocpd}): removing calendar seasonality before detection leaves the counterexample intact (phase removal keeps traffic shortest alongside the longest optimal span; seasonal differencing weakens its rank to third--fourth, still non-positive), and a changepoint-model-free drift measure independently recovers the expected negative relation ($\rho = -0.90$, Holm-adjusted $p = 0.033$ within the iTransformer statistic family). Regime persistence alone cannot explain the effective span.

\subsection{The two-axis account}

The counterexample points to the missing variable: what survives a regime change. Levels, trends, and regime parameters do not survive; periodic phase and long-memory dependence do. We therefore model the value of an old window as
\begin{equation}
V(\ell) = \underbrace{B_{\mathrm{inv}}(\ell)}_{\text{regime-invariant content}} - \underbrace{C_{\mathrm{drift}}(\ell)}_{\text{pollution cost}},
\label{eq:twoaxis}
\end{equation}
where $B_{\mathrm{inv}}$ is the information the window carries about periodic and long-memory structure (invariant across regimes), and $C_{\mathrm{drift}}$ is the cost of training on a window whose conditional mapping has expired. The EHS is the span at which the marginal benefit equals the marginal cost. This account explains the full pattern: exchange rate has high drift and weak persistent periodicity, so cost dominates at a few dozen steps; electricity and traffic switch regimes constantly but retain strong daily and weekly cycles, so distant history keeps paying; the ETT datasets sit between.

\subsection{A power law for the harm axis}
\label{sec:powerlaw}

The harm axis admits a quantitative law. A bias--variance argument predicts the optimal window under drift: bias grows with the drift accumulated over the window, variance shrinks with its length, and for continuous linear drift the predicted exponent is $W^* \propto \lambda^{-2/3}$ (derived in Appendix \ref{app:powerlaw}). We test this on controlled synthetic series whose drift rate $\lambda$ (shock probability per step) is parameterized exactly, with periodic content removed, training a small iTransformer under the equal-data-budget protocol. A first sweep over six drift tiers gives log-log slopes $-0.31$ and $-0.37$ with $R^2 = 0.82$ and $0.92$ (grid-argmin and parabolic-interpolation estimators). A hardened sweep over ten drift tiers with a refined nine-point lookback grid gives slopes $-0.29$ and $-0.32$ with $R^2 = 0.64$ and $0.76$, bootstrap intervals that contain $-\tfrac13$ and exclude $-\tfrac23$ on both estimators, and exposes a mid-drift plateau (Appendix \ref{app:powerlaw}):
\begin{equation}
W^* \propto \lambda^{-\alpha}, \qquad \alpha \approx 0.29\text{--}0.37,
\label{eq:law}
\end{equation}
because the bias of shock-driven drift scales differently from the continuous case (the shock regime sits between the naive Poisson-step prediction $-\tfrac12$ and the trend-step $-\tfrac14$). The law is family-scoped: DLinear stays short (96--480) with no monotone $\lambda$ dependence and both estimator intervals containing zero, so the power law is a coarse first description of the instance-normalized attention family, not a universal one. Two controls establish that this is the harm axis at work rather than an artifact: the drift statistic is calibrated strictly monotone in $\lambda$, and adding strong periodicity at fixed $\lambda$ extends the measured span by up to 4.3x, as the two-axis model predicts. On real benchmarks the fitted constant tracks weather (predicted 319 from the interpolation estimator, measured 336; the grid-refit constant used by the auto-decay chain in Section \ref{sec:method} gives 287) but overestimates exchange (by 1.5--2.5x) and underestimates ETTh1 (by 3x) and electricity; the residual mismatch is carried by the benefit axis, yielding the combined form $\log W^* = \max(\text{harm law}(\lambda),\, \text{benefit floor}(\text{periodicity}))$ with the constant calibrated on synthetic drift and recalibration to real drift left open (Appendix \ref{app:powerlaw}).

\subsection{Pilot predictability}
\label{sec:predict}
If the two-axis account is right, the effective span should be partially predictable from context statistics without any training. We compute per dataset, on the train split only, a drift index (variance of segment means relative to total variance), a long-memory score (autocorrelation at a lag of 25\% of the series), and a detrended spectral periodicity score, and correlate them with the measured long-context benefit $e(96) - e(2880)$. Over the six benchmark datasets of the pilot, the drift index correlates at Spearman $\rho = -0.94$ and the long-memory score at $\rho = +0.83$ ($p = 0.005$ and $0.042$, $n = 6$), while a naive spectral-peak periodicity score is uninformative ($\rho \approx 0$) because low-frequency trend power masquerades as periodicity, a measurement pitfall we document in Appendix \ref{app:stats}. A full-matrix revalidation with leave-one-dataset-out generalization strengthens this to twenty-eight samples over fourteen datasets and two model families in Section \ref{sec:predictfull}.

\subsection{Architectures differ in how they can spend history}
\label{sec:mechanism}

The model side of Table \ref{tab:bestl} shows that some families cannot exploit long context even when the data rewards it. Our measurements identify two mechanical causes.

\paragraph{Attention flattens.}
In PatchTST, attention is over patch tokens, so longer lookbacks mean more keys per query. On exchange rate, the normalized attention entropy of the second layer rises from 0.84 ($L{=}96$) to 0.91 ($L{=}2880$), approaching the uniform ceiling; attention maps become broad vertical stripes rather than concentrated columns (Figure \ref{fig:attn}). Contrary to what the attention-sink literature would suggest \citep{streamingllm}, no sink forms to absorb stale mass: the mass on the oldest patch \emph{falls} from 0.14 to 0.007. Flattening is not harmful by itself, though: the dilution is even stronger on electricity, the largest long-context beneficiary (normalized entropy $0.66 \to 0.95$), so whether it hurts depends on whether the diluted far content is noise or regime-invariant structure (Appendix \ref{app:attn}). The degradation is a loss of discrimination, not a loss of recency.

\paragraph{Free capacity to gate is not learned gating.}
A per-channel linear map could in principle assign zero weight to distant lags. Reading the trained DLinear weights on exchange rate at $L{=}2880$ shows the opposite: only 6\% of the seasonal-branch weight mass lies within the most recent 96 lags, and 73\% lies beyond lag 720 with no periodic structure, consistent with its error explosion (0.636). Empirical risk minimization does not prune stale inputs; it uses them.

\paragraph{The projection head reads what attention masks cannot hide.}
A recency-masked PatchTST variant that restricts attention to the most recent 336 steps reduces the long-context error on exchange by 34\% (0.399 $\to$ 0.263 at $L{=}2880$, closing about 45\% of the degradation gap), while a learnable null-logit variant (denominator bias) does nothing (0.399). The partial recovery, however, remains far from the short-window performance (0.095), because the flattening projection head still reads all 360 patch embeddings. Staleness must be controlled at the input, before any internal pathway can leak it (Section \ref{sec:method}). On electricity, a long-context beneficiary, the same mask does not hurt (0.127 vs.\ 0.130), consistent with the dual-channel picture: attention wants selectivity, while the projection head can safely keep the full window.
